\uuid{i8L9}
\exo7id{1987}
\titre{exo7 1987}
\auteur{ridde}
\organisation{exo7}
\datecreate{1999-11-01}
\isIndication{false}
\isCorrection{false}
\chapitre{Isométrie vectorielle}
\sousChapitre{Isométrie vectorielle}
\cpge{mpsi-mp2i.matrices-applications-lineaires.a, pcsi.matrices-applications-lineaires.a, ptsi.matrices-applications-lineaires.a, mp-mpi.endomorphismes-euclidien.b, pc.endomorphismes-euclidiens.b, psi.endomorphismes-euclidien.b, pt.isometries-euclidien.b, mp-mpi.endomorphismes-euclidien.e, psi.endomorphismes-euclidien.e, pt.isometries-euclidien.d}
\module{Géométrie}
\niveau{L2}
\difficulte{}

\contenu{
\question{
Determiner la matrice dans la base canonique de $\Rr^{3}$
du retournement d'axe $\Rr (1, 2, 1)$.
}
}
